%%%%%%%%%%%%%%%%%%%%%%%%%%%%%%%%%%%%%%%%%%%%%%%%%%%%%%%%%%%%%%%%%%%%%%%%%%%%%%%%
% LateX
%
% Author: FOSSEE 
% Copyright (c) 2009, FOSSEE, IIT Bombay
%%%%%%%%%%%%%%%%%%%%%%%%%%%%%%%%%%%%%%%%%%%%%%%%%%%%%%%%%%%%%%%%%%%%%%%%%%%%%%%%

\documentclass{beamer}
\usepackage[latin1]{inputenc}
\usepackage[T1]{fontenc}
\usepackage{fixltx2e}
\usepackage{graphicx}
\usepackage{longtable}
\usepackage{float}
\usepackage{wrapfig}
\usepackage{soul}
\usepackage{textcomp}
\usepackage{marvosym}
\usepackage{wasysym}
\usepackage{latexsym}
\usepackage{amssymb}
\usepackage{hyperref}
\tolerance=1000
\usepackage[english]{babel} \usepackage{ae,aecompl}
\usepackage{mathpazo,courier,euler} \usepackage[scaled=.95]{helvet}
\usepackage{listings}
\lstset{
  language=TeX,
  basicstyle=\ttfamily\bfseries,
  commentstyle=\ttfamily\color{blue},
  stringstyle=\ttfamily\color{orange},
  showstringspaces=false,
  breaklines=true,
  postbreak = \space\dots
}

\newcommand{\typ}[1]{\lstinline{#1}}

\mode<presentation>
{
  \usetheme{Warsaw}
  \useoutertheme{infolines}
  \setbeamercovered{transparent}
}


\setbeamercolor{emphbar}{bg=blue!20, fg=black}
\newcommand{\emphbar}[1]


\begin{document}

\begin{frame}

\begin{center}
\vspace{12pt}
\textcolor{blue}{\huge {\LaTeX}:Tables, Figures and Floats}
\end{center}
\vspace{18pt}
\begin{center}
\vspace{10pt}
\includegraphics[scale=0.95]{../images/fossee-logo.png}\\
\vspace{5pt}
\scriptsize Developed by FOSSEE Team, IIT-Bombay. \\ 
\scriptsize Funded by National Mission on Education through ICT\\
\scriptsize  MHRD,Govt. of India\\
\includegraphics[scale=0.15]{../images/iitb-logo.jpg}\\
\end{center}
\end{frame}


\begin{frame}
  \frametitle{Objectives}
  At the end of this session, you will be able to:
  \begin{itemize}
  \item Learn how to add figures to your document
  \item Include tabular environments
  \end{itemize}
\end{frame}


\begin{frame}
  \frametitle{Prerequisite}
  Spoken Tutorial on:
  \begin{enumerate}
  \item {\LaTeX} - Getting Started
  \item {\LaTeX} - Introduction
  \item {\LaTeX} - Structuring the Content
  \item {\LaTeX} - Typesetting Text
  \end{enumerate}
\end{frame}

\begin{frame}[fragile]
  \frametitle{Figures}
  \begin{itemize}
  \item The \typ{graphicx} package allows us to insert graphics
  \item \lstinline+\usepackage{graphicx}+
  \item To add a graphic, use \lstinline{\includegraphics} command
  \item Use relative path to the image
  \end{itemize}
  \tiny See rev17 of \typ{hg}
\end{frame}

\begin{frame}[fragile]
  \frametitle{\lstinline{includgraphics}}
  It takes following optional arguments
  \begin{itemize}
  \item \lstinline+scale+ --- specifies the factor by which to scale
    the image 
  \item \lstinline+height+, \lstinline+width+ --- If only one of them
    is specified, aspect ratio is maintained 
  \item \lstinline+keepaspectratio+ --- boolean value to keep aspect
    ratio or not 
  \item \lstinline+angle+ --- specify by what angle the image should
    be rotated 
  \end{itemize}
\end{frame}

\begin{frame}[fragile]
  \frametitle{Floats}
  \begin{itemize}
  \item Graphics (\& Tables) are special because they cannot be broken
    across pages 
  \item They are ``floated'' to the next page, if they don't fit in
    the current page 
  \item Enclose graphic within \lstinline+figure+ environment to make
    it float 
  \item Figure environment takes additional parameter for location of
    float 
  \end{itemize}
  \begin{table}
    \caption{Permission Specifiers}
    
    \begin{tabular}{|c|c|}
      Specifier & Permission\\\hline
      t & Top of page\\
      b & Bottom of page\\
      p & Separate page for floats\\
      h & here (the same place where command appears in source)\\
      ! & override \LaTeX's internal parameters for good position
    \end{tabular}
  \end{table}
\end{frame}

\begin{frame}
  \frametitle{Captions and References}
  \begin{itemize}
  \item Figure environment allows us add a caption
  \item To place the image in the center we enclose it in the
    \lstinline+center+ environment 
  \item We can label images too
  \item label should be added after the caption command
  \item Figures are auto numbered
  \end{itemize}
  \tiny See rev17 of \typ{hg}
\end{frame}

\begin{frame}[frame]
  \frametitle{Tables}
  \begin{itemize}
  \item \lstinline+tabular+ is used to typeset a table
  \item It is enclosed in a \lstinline+table+ environment to make it a
    float 
  \item \lstinline+table+ environment also gives captions, auto
    numbering  
  \end{itemize}
\end{frame}


\begin{frame}[fragile]
  \frametitle{\lstinline+tabular+}
  \begin{itemize}
  \item tabular takes formatting of each column as argument
  \end{itemize}

  \begin{table}
    \caption{tabular environment}
    
    \begin{tabular}{|l|l|}
      \lstinline+l+ & left justified column content\\\hline
      \lstinline+r+ & right justified column content\\\hline
      \lstinline+c+ & centered column content\\\hline
      \lstinline+|+ & produces a vertical line\\
    \end{tabular}
  \end{table}
  \begin{itemize}
  \item also takes an optional parameter for specifying position of
    table 
  \item \lstinline+t+ for top, \lstinline+b+ for bottom, \lstinline+c+
    for center 
  \item each column of table is separated by \&
  \item each row is separated by newline \lstinline{\\}
  \item \lstinline+\hline+ give a horizontal line between two rows
  \end{itemize}
  \tiny See rev18 of \typ{hg}
\end{frame}

\begin{frame}[fragile]
  \frametitle{List of Tables, Figures}
  \begin{itemize}
  \item \lstinline+\listoftables+ -- to add a list of tables
  \item \lstinline+\listoffigures+ -- to add a list of figures
  \end{itemize}
\end{frame}

\begin{frame}[fragile]
	\frametitle{Summary...}
	\begin{itemize}
	\item Add graphics to a LateX document
	\item Include tabular environments in a LateX document
	\end{itemize}
\end{frame}

\begin{frame}[fragile]
\frametitle{Evaluation}
\begin{enumerate}
\item Why the figures and tables are called as floats?
\item What are the input parameters the figure environment take to make it float?
\item What is the mandatory arguement in tabular environment specification?
\end{enumerate}
\end{frame}
\begin{frame}

\frametitle{Solutions}
\begin{enumerate}
\vspace{15pt}
\item Figures and tables are called as floats because they float to the nearest location.
\item The following can be used as parameters to a figure environment to make it float.
\begin{itemize}
\item t - Top of the page.
\item b - Bottom of the page.
\item p - Separate pages for float
\item h - Here(the place where command appears)
\item !- To override Latex internal parameters for good position
\end{itemize}

\item Alignment of each column
\end{enumerate}
\end{frame}

\begin{frame}

\begin{block}{}
  \begin{center}
  \textcolor{blue}{\Large THANK YOU!} 
  \end{center}
  \end{block}
\begin{block}{}
  \begin{center}
    For more Information, visit our website\\
    \url{http://fossee.in/}
  \end{center}  
  \end{block}
\end{frame}

\end{document}

